\documentclass{beamer}
\usepackage[utf8]{inputenc}
\usepackage[spanish]{babel}
\usepackage{listings}
\usepackage{xcolor}
\usepackage{tikz}
\usepackage{tcolorbox}
\usepackage{fontawesome5}

% Configuración de colores - Gruvbox Light con marrón principal
\definecolor{bg0}{HTML}{fbf1c7}
\definecolor{bg1}{HTML}{ebdbb2}
\definecolor{bg2}{HTML}{d5c4a1}
\definecolor{bg3}{HTML}{bdae93}
\definecolor{fg0}{HTML}{282828}
\definecolor{fg1}{HTML}{3c3836}
\definecolor{fg2}{HTML}{504945}
\definecolor{aqua}{HTML}{689d6a}
\definecolor{aquabright}{HTML}{427b58}
\definecolor{blue}{HTML}{458588}
\definecolor{red}{HTML}{cc241d}
\definecolor{orange}{HTML}{d65d0e}
\definecolor{yellow}{HTML}{d79921}
\definecolor{green}{HTML}{98971a}
\definecolor{purple}{HTML}{b16286}
\definecolor{brown}{HTML}{AB886D}
\definecolor{brownbright}{HTML}{8B6849}

% Colores para código
\definecolor{codegreen}{HTML}{98971a}
\definecolor{codegray}{HTML}{7c6f64}
\definecolor{codepurple}{HTML}{b16286}
\definecolor{backcolour}{HTML}{f9f5d7}

% Configuración de código
\lstdefinestyle{pythonstyle}{
    backgroundcolor=\color{backcolour},   
    commentstyle=\color{codegreen},
    keywordstyle=\color{brownbright}\bfseries,
    numberstyle=\tiny\color{codegray},
    stringstyle=\color{codepurple},
    basicstyle=\ttfamily\footnotesize\color{fg0},
    breakatwhitespace=false,         
    breaklines=true,                 
    captionpos=b,                    
    keepspaces=true,                 
    numbers=left,                    
    numbersep=3pt,                  
    showspaces=false,                
    showstringspaces=false,
    showtabs=false,                  
    tabsize=2,
    language=Python,
    frame=single,
    rulecolor=\color{brown},
    xleftmargin=8pt,
    xrightmargin=8pt
}

\lstset{style=pythonstyle}

% Tema personalizado con marrón principal
\usetheme{Madrid}
\setbeamercolor{structure}{fg=brown}
\setbeamercolor{palette primary}{bg=brown,fg=bg0}
\setbeamercolor{palette secondary}{bg=brownbright,fg=bg0}
\setbeamercolor{palette tertiary}{bg=fg1,fg=bg0}
\setbeamercolor{palette quaternary}{bg=brown,fg=bg0}
\setbeamercolor{block title}{bg=brown,fg=bg0}
\setbeamercolor{block body}{bg=bg1,fg=fg0}
\setbeamercolor{frametitle}{bg=brown,fg=bg0}
\setbeamercolor{title}{fg=fg0}
\setbeamercolor{subtitle}{fg=fg1}
\setbeamercolor{normal text}{fg=fg0,bg=bg0}
\setbeamercolor{item}{fg=brown}
\setbeamercolor{section in toc}{fg=brown}

% Cajas personalizadas con marrón
\newtcolorbox{conceptbox}[1]{
    colback=bg1,
    colframe=brown,
    fonttitle=\bfseries\small,
    title=#1,
    coltitle=fg0,
    left=3pt,
    right=3pt,
    top=3pt,
    bottom=3pt,
    boxsep=2pt
}

\newtcolorbox{notebox}{
    colback=yellow!15,
    colframe=orange,
    leftrule=2mm,
    rightrule=0mm,
    toprule=0mm,
    bottomrule=0mm,
    arc=0mm,
    left=3pt,
    right=3pt,
    top=2pt,
    bottom=2pt
}

\title{\textbf{Introducción a la Programación con Python}}
\subtitle{Semana 2: Bucles y Estructuras de Repetición}
\author{$\nabla$Nablabs}
\date{}
\begin{document}

\begin{frame}
\titlepage
\end{frame}

\begin{frame}
\frametitle{Contenido}
\tableofcontents
\end{frame}

\section{¿Qué es un bucle?}

\begin{frame}{Bucles: Repetir acciones \faRedo}
\begin{conceptbox}{¿Para qué sirven los bucles?}
\small
\begin{itemize}
    \item[\faSync] Permiten ejecutar un bloque de código varias veces
    \item[\faSync] Evitan la repetición manual de instrucciones
    \item[\faSync] Son esenciales para procesar datos y automatizar tareas
\end{itemize}
\end{conceptbox}

\vspace{0.2cm}
\begin{center}
\large\textcolor{brown}{\faCogs\ \texttt{while}, \texttt{for}}
\end{center}
\end{frame}

\section{Bucle while}

\begin{frame}[fragile]{Bucle while \faSync}
\begin{conceptbox}{¿Cómo funciona?}
Repite un bloque de código mientras una condición sea verdadera
\end{conceptbox}

\begin{lstlisting}
contador = 1
while contador <= 5:
    print(contador)
    contador += 1
\end{lstlisting}

\vspace{-0.2cm}
\begin{tcolorbox}[colback=green!10,colframe=green,title={\small\faArrowCircleRight \ Salida},coltitle=fg0,left=3pt,right=3pt,top=2pt,bottom=2pt,boxsep=1pt]
\small\texttt{1\\
2\\
3\\
4\\
5}
\end{tcolorbox}

\vspace{-0.1cm}
\begin{notebox}
\small\faInfoCircle\ El bucle puede ejecutarse 0 o más veces
\end{notebox}
\end{frame}

\section{Bucle for}

\begin{frame}[fragile]{Bucle for \faRepeat}
\begin{conceptbox}{¿Cómo funciona?}
Itera sobre los elementos de una secuencia (lista, string, etc.)
\end{conceptbox}

\begin{lstlisting}
frutas = ["manzana", "banana", "cereza"]
for fruta in frutas:
    print(fruta)
\end{lstlisting}

\vspace{-0.2cm}
\begin{tcolorbox}[colback=green!10,colframe=green,title={\small\faArrowCircleRight \ Salida},coltitle=fg0,left=3pt,right=3pt,top=2pt,bottom=2pt,boxsep=1pt]
\small\texttt{manzana\\
banana\\
cereza}
\end{tcolorbox}

\vspace{-0.1cm}
\begin{notebox}
\small\faLightbulb\ El bucle for recorre cada elemento automáticamente
\end{notebox}
\end{frame}

\section{Iteración sobre secuencias}

\begin{frame}[fragile]{Iterar sobre secuencias \faListOl}
\begin{columns}[T]
\column{0.48\textwidth}
\begin{lstlisting}
texto = "Python"
for letra in texto:
    print(letra)
\end{lstlisting}

\vspace{-0.2cm}
\begin{tcolorbox}[colback=green!10,colframe=green,title={\small\faArrowCircleRight \ Salida},coltitle=fg0,left=3pt,right=3pt,top=2pt,bottom=2pt,boxsep=1pt]
\small\texttt{P\\
y\\
t\\
h\\
o\\
n}
\end{tcolorbox}

\column{0.48\textwidth}
\begin{lstlisting}
numeros = [1, 2, 3, 4]
for n in numeros:
    print(n * 2)
\end{lstlisting}

\vspace{-0.2cm}
\begin{tcolorbox}[colback=green!10,colframe=green,title={\small\faArrowCircleRight \ Salida},coltitle=fg0,left=3pt,right=3pt,top=2pt,bottom=2pt,boxsep=1pt]
\small\texttt{2\\
4\\
6\\
8}
\end{tcolorbox}
\end{columns}
\end{frame}

\section{Función range()}

\begin{frame}[fragile]{Función range() \faRuler}
\begin{columns}[T]
\column{0.48\textwidth}
\begin{lstlisting}
for i in range(5):
    print(i)
\end{lstlisting}
\vspace{-0.2cm}
\begin{tcolorbox}[colback=green!10,colframe=green,title={\small\faArrowCircleRight \ Salida},coltitle=fg0,left=3pt,right=3pt,top=2pt,bottom=2pt,boxsep=1pt]
\small\texttt{0\\
1\\
2\\
3\\
4}
\end{tcolorbox}

\vspace{0.2cm}
\begin{lstlisting}
for i in range(2, 7):
    print(i)
\end{lstlisting}
\vspace{-0.2cm}
\begin{tcolorbox}[colback=green!10,colframe=green,title={\small\faArrowCircleRight \ Salida},coltitle=fg0,left=3pt,right=3pt,top=2pt,bottom=2pt,boxsep=1pt]
\small\texttt{2\\
3\\
4\\
5\\
6}
\end{tcolorbox}
\column{0.48\textwidth}
\begin{lstlisting}
for i in range(1, 10, 2):
    print(i)
\end{lstlisting}
\vspace{-0.2cm}
\begin{tcolorbox}[colback=green!10,colframe=green,title={\small\faArrowCircleRight \ Salida},coltitle=fg0,left=3pt,right=3pt,top=2pt,bottom=2pt,boxsep=1pt]
\small\texttt{1\\
3\\
5\\
7\\
9}
\end{tcolorbox}
\end{columns}
\end{frame}

\section{Sentencias break y continue}

\begin{frame}[fragile]{break y continue \faStepForward}
\begin{columns}[T]
\column{0.48\textwidth}
\begin{lstlisting}
for i in range(1, 10):
    if i == 5:
        break
    print(i)
\end{lstlisting}
\vspace{-0.2cm}
\begin{tcolorbox}[colback=green!10,colframe=green,title={\small\faArrowCircleRight \ Salida},coltitle=fg0,left=3pt,right=3pt,top=2pt,bottom=2pt,boxsep=1pt]
\small\texttt{1\\
2\\
3\\
4}
\end{tcolorbox}

\column{0.48\textwidth}
\begin{lstlisting}
for i in range(1, 6):
    if i == 3:
        continue
    print(i)
\end{lstlisting}
\vspace{-0.2cm}
\begin{tcolorbox}[colback=green!10,colframe=green,title={\small\faArrowCircleRight \ Salida},coltitle=fg0,left=3pt,right=3pt,top=2pt,bottom=2pt,boxsep=1pt]
\small\texttt{1\\
2\\
4\\
5}
\end{tcolorbox}
\end{columns}
\end{frame}

\section{Bucles infinitos y cómo evitarlos}

\begin{frame}[fragile]{Bucles infinitos \faInfinity}
\begin{conceptbox}{¿Qué es un bucle infinito?}
Un bucle que nunca termina porque la condición siempre es verdadera
\end{conceptbox}

\begin{lstlisting}
# ¡Cuidado! Esto nunca termina
while True:
    print("Bucle infinito")
    # Debe haber un break o condición de salida
\end{lstlisting}

\vspace{-0.2cm}
\begin{tcolorbox}[colback=yellow!10,colframe=orange,title={\small\faExclamationTriangle \ Salida (parcial)},coltitle=fg0,left=3pt,right=3pt,top=2pt,bottom=2pt,boxsep=1pt]
\small\texttt{Bucle infinito\\
Bucle infinito\\
Bucle infinito\\
...}
\end{tcolorbox}

\vspace{-0.1cm}
\begin{notebox}
\small\faExclamationTriangle\ Siempre asegúrate de que la condición cambie en algún momento
\end{notebox}
\end{frame}

\section{Validación de entrada con bucles}

\begin{frame}[fragile]{Validación de entrada \faCheckCircle}
\begin{lstlisting}
while True:
    edad = input("Introduce tu edad: ")
    if edad.isdigit():
        edad = int(edad)
        break
    print("Por favor, ingresa un número válido.")
\end{lstlisting}

\vspace{0.2cm}
\begin{columns}[T]
\column{0.48\textwidth}
\begin{tcolorbox}[colback=blue!10,colframe=blue,title={\small\faArrowDown \ Entrada},coltitle=fg0,left=2pt,right=2pt,top=2pt,bottom=2pt]
\small\texttt{Introduce tu edad: hola\\
Por favor, ingresa un número válido.\\
Introduce tu edad: 25}
\end{tcolorbox}

\column{0.48\textwidth}
\begin{tcolorbox}[colback=green!10,colframe=green,title={\small\faArrowUp \ Salida final},coltitle=fg0,left=2pt,right=2pt,top=2pt,bottom=2pt]
\small\texttt{edad = 25}
\end{tcolorbox}
\end{columns}
\end{frame}

\section{Ejemplo integrador}

\begin{frame}[fragile]{Ejemplo práctico \faUserCog}
\begin{lstlisting}
# Sumar los números del 1 al 100
suma = 0
for i in range(1, 101):
    suma += i
print(f"La suma es: {suma}")
\end{lstlisting}

\vspace{-0.2cm}
\begin{tcolorbox}[colback=green!10,colframe=green,title={\small\faArrowCircleRight \ Salida},coltitle=fg0,left=3pt,right=3pt,top=2pt,bottom=2pt,boxsep=1pt]
\small\texttt{La suma es: 5050}
\end{tcolorbox}
\end{frame}

\begin{frame}{Resumen \faListUl}
\begin{conceptbox}{¿Qué aprendimos hoy?}
\small
\begin{itemize}
    \item[\faCheck] Bucle while
    \item[\faCheck] Bucle for
    \item[\faCheck] Iteración sobre secuencias
    \item[\faCheck] Función range()
    \item[\faCheck] break y continue
    \item[\faCheck] Evitar bucles infinitos
    \item[\faCheck] Validación de entrada
\end{itemize}
\end{conceptbox}

\vspace{0.3cm}
\begin{center}
\large\textcolor{brown}{\faCode\ ¡Practica con bucles!}
\end{center}
\end{frame}

\begin{frame}
\begin{center}
\Huge\textcolor{brown}{\faQuestionCircle}\\
\vspace{0.5cm}
\Huge ¡Preguntas!\\
\vspace{1cm}
\Large\textcolor{brownbright}{Gracias por su atención}
\end{center}
\end{frame}

\end{document}
